This is a small study. Five hand-built 6x9 mazes, one fixed change time per task, one set of hyperparameters per agent, 20 seeds for the main table and 10 for the sweeps. Per-seed standard deviations are often a large fraction of the mean (many seeds have a handful of goals), so differences below roughly a quarter of the mean are not resolved; $p$-values are uncorrected over many comparisons and should be read as descriptive. The baselines are our reimplementations from the original descriptions and textbook, not reference code, and details (tie-breaking, how untried actions enter the model, the priority threshold) can matter. We tuned nothing, which is fair across methods but means each may be far from its best: Dyna-Q+'s collapse at larger $\kappa$ shows how strongly this matters, and $\epsilon$ and $\alpha$ interact with $n$. We swept $n$ only for Dyna-Q, Dyna-Q+ and PS and with $n=1$ the early learning is so slow that the staleness effect in H2 is confounded, so H2's first half is inconclusive rather than disproved at larger scale. The sweeps use 10 seeds, which leaves the stochastic-maze comparison with Q-learning underpowered. The model stores only the last outcome; a count-based or sampling stochastic model, which might remove the slip decline, is out of scope and untested, so we do not know how much of the decline is due to the last-outcome model rather than to planning itself (there is no correct-model control). We did not log per-episode path choices, so mechanism statements (e.g.\ Dyna-Q never finds the shortcut) are inferences from reward counts. A full study would add more mazes and change times, larger grids, tuned hyperparameters with a held-out protocol, interval estimators from the literature~\cite{agarwal2021deep}, and function approximation.
